% =============================================================================
% Draft for report/report.tex sec:method (lines 80-91 of the current skeleton).
% Part of the P3 report-writing effort -- issue #70 (p3.1-method).
% ~650 words. Covers: channel taxonomy, claim types, degradation conditions
% and verdicts, verifiability tiers, coding procedure, reliability protocol
% (including the D2 failure and fix), and the R-01-R-07 operating constraints.
% =============================================================================

\section{Method}
\label{sec:method}

Every public claim entering the ledger (\texttt{ledger/claims.csv}) is coded on four
independent axes against a frozen codebook (\texttt{docs/codebook.md}, v1.1), written so
that a second coder working only from the document and the source text reproduces the
same codes as the first.

\paragraph{Channel taxonomy and claim types.} Each claim is assigned a primary
evidentiary channel plus zero or more corroborating channels, drawn from a seven-value
taxonomy ordered by fragility and agent-writability: C1 (raw chain-of-thought), C2
(inter-agent messages), C3 (tool-call/action traces), C4 (defender-side telemetry), C5
(human testimony), C6 (lab assertion with no stated substrate), and UNSTATED (no basis
given at all). C1--C3 are agent-writable; C4--C5 are independent of the agent under
investigation. Each claim is also typed by the kind of assertion it makes: T1 (event),
T2 (mechanism), T3 (intent), T4 (counterfactual), or T5 (assurance). Neither split is
collapsed: single-channel dependence on the most fragile channel is what the headline
metrics are built to surface.

\paragraph{Degradation conditions and verdicts.} Every claim is stress-tested against
three conditions. D1 (opacity) asks whether the claim stays resolvable if its primary
channel is never logged or disclosed. D2 (unfaithfulness) asks whether the claim's
warrant survives if the channel's content, though available, is not a reliable report of
the computation that produced it. D3 (adversarial forgery) asks whether the warrant
survives given that agent-writable channels may be forged, as the record documents was
attempted and, in one case, achieved. Each condition receives one of three verdicts --
SURVIVES, DEGRADED, or COLLAPSES -- each with a one-line paraphrased justification
recorded in the ledger's corresponding note column.

\paragraph{Outside-verifiability.} Independent of channel and degradation coding, each
claim is assigned one of three verifiability tiers: V1 (checkable by any third party from
public artifacts), V2 (checkable only with lab or platform cooperation), or V3 (not
checkable by any stated procedure). V3 is assigned by definition to any claim -- typically
a T5 assurance claim -- that states no verification procedure at all; "not checkable" is
a value in an exhaustive partition, never a blank or N/A.

\paragraph{Coding procedure.} Claims are paraphrased into the ledger, never quoted
verbatim, for copyright reasons. Each row carries the source organization, URL, and a
retrieval date, so sourcing is dated rather than assumed stable. A single coder applied
the frozen codebook, with edge cases resolved by explicit, priority-ordered rules (e.g., a
fixed fragility ranking breaks primary/corroborating ties) rather than case-by-case
judgment, so a second coder is expected to land on the same codes.

\paragraph{Reliability protocol.} A 25-row subsample was stripped of its verdicts and
re-coded cold by a second pass with no access to the original codes, then compared
against the original via Cohen's $\kappa$ per condition against a pre-registered 0.6 bar.
D1 ($\kappa=0.7788$) and D3 ($\kappa=0.7794$) passed comfortably. D2 did not
($\kappa=0.5690$), despite 92\% raw agreement, because the D2 verdict distribution's skew
toward SURVIVES inflates chance agreement. Diagnosis
(\texttt{adr/0001-d2-self-report-corroboration.md}) found the codebook let coders reuse
D1/D3's corroboration-rescues-the-verdict logic for D2, when unfaithfulness can be a
systematic narration habit rather than a one-off gap: a channel merely repeating the same
kind of self-reported content, or covering a mismatched scope, does not rescue a D2
verdict the way it rescues D1 or D3. The fix was logged as a dated ADR, and the ledger's
entire D2 column was recoded against the corrected rule before the freeze, while D1 and D3
verdicts were left untouched. This failure-and-fix is reported as a strength, not
softened: it is the clearest evidence in the project that the reliability check was a real
test, not a formality, catching a genuine rule gap rather than coder inattention, with the
fix fully auditable in the diff between codebook versions.

\paragraph{Operating constraints.} The instrument runs under seven numbered requirement
rules (R-01 through R-07) that keep the coding auditable and replicable, covering source
retrieval-dating, channel-taxonomy assignment, script-computed metrics, blinded
reliability reporting, MRFM clause verifiability, no reproducible tamper or spoof
technique, and dated-ADR logging of any post-freeze codebook change.
